\section{Finite-precision gates and stopping-time guarantees}
\label{sec:statistics}
The descriptor separation controls class and sign errors. It does not
make measured arclength origins exact. We now account for that second
error, as well as the nonuniformity of laboratory-square preparation,
before applying the deterministic certificate.

\subsection{Registration error enters through cell boundaries}
Fix symmetric square grids on a rectangle whose closure lies in the
protected return rectangle, with a further collar for interpolation.
Suppose the probe reports endpoints $\widehat z=(\widehat s,\widehat t)$
whose correctly oriented exact coordinates $z$ satisfy
$\|\widehat z-z\|_\infty\le\nu\le h/4$. Class and sign identification
are exact by Proposition~\ref{prop:registry}. For a square cell $B$ of
side $h$, a changed bin assignment is possible only if $z$ lies in
$B^{+\nu}\setminus B^{-\nu}$. The area of this band is at most
$8h\nu+4\nu^2$, and its quotient probability is at most $C\nu h$
under the bounded physical density. Under laboratory-square preparation
it is at most $C\nu h+\tau_L$.

Let $J$ be a finite number of records chosen before their fresh validation
samples, let $N$ be the sample size at each of their two times, and put
\[
 r_N=\frac{\log(C(J\vee1)N/\alpha)}{N}.
\]
For sufficiently large $N$ and the bandwidths below, simultaneous cell
errors are, with conditional probability at least $1-\alpha$,
\begin{equation}\label{eq:cell-error}
 C\{h\sqrt{r_N}+r_N+\tau_L+\nu h\}.
\end{equation}
Here is a justification which allows arbitrary bounded numerical probe
errors. Apply Bernstein's inequality to the true cell indicators and to
the indicators of their boundary bands, conditional on the fixed records
and times. Their probabilities are bounded by $Ch^2+\tau_L$ because
$\nu\le h/4$. The empirical difference between a measured and an exact
cell indicator is bounded by the empirical band indicator, regardless
of how the numerical errors were chosen. Add the square-preparation bias
and use $\sqrt{\tau_Lr_N}\le(\tau_L+r_N)/2$. A union bound over
$O(Jh^{-2})$ cells proves \eqref{eq:cell-error}. For the bandwidths used
below $h^{-2}\le CN$, so the displayed logarithm suffices. This does
not assume that adversarial numerical rounding is independent across trials.

\begin{lemma}[Two-derivative acquisition with local registration error]
\label{lem:histogram}
For physical densities bounded in $C^{2,\beta}$ on the larger active
rectangle, a reconstruction from the finite cell probabilities satisfies
\begin{equation}\label{eq:density-error}
 \|\widehat f-f\|_{C^2}\le
 C\{h^\beta+\sqrt{r_N}h^{-3}+(r_N+\tau_L)h^{-4}
                                      +\nu h^{-3}\}
\end{equation}
on the simultaneous event above. The constants are uniform over genuine
record lists and their repeated entries.
\end{lemma}
\begin{proof}
The averages of $1,z,z^2$ on three unit intervals centered at
$i=-1,0,1$ form the invertible matrix with rows $(1,i,i^2+1/12)$.
Its tensor square reconstructs coordinatewise quadratic polynomials
from a $3\times3$ cell block. A degree-two Taylor expansion gives
local derivative error $O(h^{2+\beta-j})$ through order $j\le2$.
A cell-mass error $\epsilon$ becomes average-density error
$\epsilon h^{-2}$, and $j$ derivatives cost $h^{-j}$.
Patch the local polynomials with a smooth uniformly locally finite
partition of unity. Its derivative costs are offset by the corresponding
lower-order Taylor errors. Thus the $C^2$ error is bounded by
$C(h^\beta+\epsilon h^{-4})$. Inserting \eqref{eq:cell-error} proves
\eqref{eq:density-error}. The observation collar supplies all boundary
stencils. No $C^2$ estimate is deduced directly from total variation.
\end{proof}

In particular take, for small $r_N$,
\begin{equation}\label{eq:balanced}
 h=r_N^{1/(2\beta+6)},\qquad
 L\ge C h^{-(\beta+4)},\qquad
 \nu\le c h^{\beta+3}.
\end{equation}
The square bias obeys $\tau_L\le Ch^{\beta+4}$ by \eqref{eq:square}.
Each term in \eqref{eq:density-error} is then at most
$C r_N^{\beta/(2\beta+6)}$. Notice the different amplification of the
two systematic errors: square bias costs $h^{-4}$, whereas a bounded
coordinate displacement costs $h^{-3}$. Treating a shifted contact origin
as an exact intrinsic mark would miss this term.

The derivative order $K_*$ for class identification remains fixed.
Only the coordinate accuracy $\nu$ is refined. Reduce the probe accuracy
also below $\chi/16$ to keep classification exact. The identities
$h^{\beta+3}=r_N^{1/2}$ and \eqref{eq:balanced} will be useful in the
resource accounting.

\subsection{An explicit sequential procedure}
Epochs are indexed by $k=1,2,\ldots$. First make one fresh discovery
proposal as in Section~\ref{sec:acquisition}, using its fixed discovery
square and time grid. Retain any genuine descriptor, allowing repeated
entries; the list length satisfies $J_k\le k$. If the list is empty,
continue without certification. Otherwise, conditional on the whole
proposal history, validate every retained entry at its two active times
with fresh free preparations. Use
\begin{equation}\label{eq:schedule}
 N_k=(k+1)^2,\qquad
 \alpha_k=\frac{\delta}{8(k+1)^a},\qquad a\ge6,
\end{equation}
and the bandwidth, square and coordinate precision in \eqref{eq:balanced}.
For the finitely many early epochs outside its small-error range one
may record data but issue no certificate. In particular the fixed
fingerprint and protected-rectangle conditions are checked before use.

Fit physical density pairs, identify visible types and apply the rational
and defect test of Theorem~\ref{thm:certificate}. For a prescribed
accuracy $\xi$, issue a certificate only when the deterministic error
bound $Ce_k^\theta$ is at most $\xi$ and the theorem's arithmetic and
defect thresholds hold. The smallest such passing epoch is $T_\xi$.
The true local laws have common density error at most
\begin{equation}\label{eq:epoch-error}
 e_k=C\left\{\frac{\log(C(k+1)^{a+3}/\delta)}{(k+1)^2}
                              \right\}^{\beta/(2\beta+6)}
\end{equation}
except on an event of conditional probability at most $\alpha_k$.
This bound tends to zero. The current list is fixed when the fresh
validation begins; it need not have been fixed before the experiment.

\begin{proof}[Proof of Theorem~\ref{thm:main}, part (iii)]
Let $E_k$ be the accuracy event for all records validated at epoch $k$.
Conditional concentration followed by iterated conditioning gives
$\Prob(E_k^c)\le\alpha_k$. Since epochs start at one,
$\sum_{k\ge1}\alpha_k<\delta/8$. On their simultaneous intersection,
Theorem~\ref{thm:certificate} applies to every component and every
chosen deletion of every current list. Thus all issued certificates,
including at data-dependent stopping times, are correct. No union bound
over selected graphs or stopping rules is needed.

There is a finite deterministic $k_\xi$ after which
\eqref{eq:epoch-error} satisfies every required threshold. If the
saturated witness has appeared by epoch $k\ge k_\xi$ and $E_k$
holds, some complete component passes and stopping has occurred by $k$.
Therefore
\[
 \{T_\xi>k\}\subset\{K_{\rm cov}>k\}\cup E_k^c.
\]
Equations \eqref{eq:coupon} and \eqref{eq:schedule} give
\eqref{eq:main-tail}, with the chosen $a$. In particular the tail tends
to zero and $T_\xi$ is almost surely finite. Soundness used only genuine
local gates; the positive hazard was used for termination and cost, not
for the validity of an issued certificate.
\end{proof}

\subsection{Separate resource accounts}
Through epoch $m$, accepted free preparations, including one proposal
per epoch, are at most
\begin{equation}\label{eq:launch-cost}
 m+2\sum_{k=1}^m J_kN_k\le C(m+1)^4.
\end{equation}
At every square size the free probability is at least $\pi_0$ by
\eqref{eq:square}. A negative-binomial bound, or a binomial Chernoff
bound at a fixed trial budget, gives at most
$C_{\pi_0}\{n+\log(1/\gamma)\}$ attempted positions for $n$ accepted
launches with probability at least $1-\gamma$. Allocate another
$\delta/[8(k+1)^a]$ to each epoch and sum. The high-probability launch
budget through $m$ is
$O(m^4+m\log(m/\delta))$. All solid rejections, return failures and
out-of-rectangle endpoints are charged.

The largest laboratory-square half-width through $m$ is bounded by
$C(m+1)^{(\beta+4)/(\beta+3)}$; the fixed discovery square is included
in the constant. Each trace has time horizon bounded by a prior constant,
and at most $1+T_{\max}/d_0$ collisions. The local-query work obeys
\begin{equation}\label{eq:probe-cost}
 C\sum_{k=1}^m\{1+J_kN_k\}
 P_{K_*}\!\left(c\min\{\chi,\nu_0,r_k^{1/2}\}\right),\qquad
 r_k=\frac{\log(C(k+1)^{a+3}/\delta)}{(k+1)^2}.
\end{equation}
This is a separate account, not an omitted factor in the preparation
count. If $P_{K_*}(\nu)\le C\nu^{-s}$, it is $O(m^{4+s})$.

Before epoch $k$, the event that its work is needed is determined by the
past. Its conditional expected launch cost is at most $Ck^3$ by the
uniform free-position bound. Multiplying this by the tail at epoch $k-1$
and summing proves finite expected launch cost when $a>4$, in particular
for $a=6$. Under the polynomial probe-cost assumption, the corresponding
term is $Ck^{3+s}$; its sum is finite when $a>4+s$. One can select this
larger error-spending exponent in advance without changing the powers
in \eqref{eq:balanced}. If apparatus travel per operation is also
bounded by a power of $L_k$, its additional exponent must likewise enter
this summability condition. No expected cost is claimed for arbitrary
unbounded probe cost or unmodelled travel and setup.

Compact analytic fitting and continuation are still existence procedures,
not algorithms with a proved polynomial bit complexity. The finite
integer recovery is effective after its separated inputs are obtained.
These distinctions also apply to every headline complexity statement.
The retained pointwise margin-exhaustion theorem has no uniform hazard
or uniform expected-cost assertion over its exhaustion levels.
